\begin{frame}{프로젝트 구조}
  \begin{itemize}
    \item \kb{데이터} --- Kaggle, UCI Machine Learning Repository
      \textbf{중에서 정형(행/열) 데이터 하나} --- 이미지·텍스트보다
      표 형태가 Week 2--7 모델 검증에 적합.
    \item \kb{적용} --- Week 2--7 중 \textbf{최소 두 가지 모델}
      (예: 로지스틱회귀 + GBDT)을 같은 데이터에 적용해 비교.
    \item \kb{필수: 수식적 정당화 $\geq$1개} ---
      \begin{enumerate}
        \item 정규화 적용했다면 Week 6 편향-분산 관점에서 왜 성능이
          바뀌었는지,
        \item GBDT 썼다면 Week 7 SHAP으로 예측 하나를 분해,
        \item 모델 여러 개 비교했다면 Week 2.3 PR-AUC로 불균형을
          고려한 비교.
      \end{enumerate}
    \item \kb{검증 원칙} --- Week 6.3의 train/validation/test 3분할
      준수, \textbf{테스트 성능은 딱 한 번}만 확인.
  \end{itemize}
\end{frame}
